\documentclass{article}
\title{ECE 311 -- Homework 5}
\usepackage[margin=1in]{geometry}
\usepackage{graphicx}
\usepackage{pdfpages}
\usepackage{fancyhdr}
\usepackage{enumitem}
\usepackage{pdfpages}
\usepackage{amsmath}
\usepackage{amssymb}
\usepackage{float}
\usepackage{xcolor}
\pagestyle{fancy}
\fancyhf{}
\fancyhead[L]{Gianni Avilan-Losee}
\fancyhead[C]{ECE 311}
\fancyhead[R]{Homework 5}
\fancyfoot[C]{\thepage}
\usepackage{microtype}
\usepackage{textcomp, gensymb}
\author{Gianni Avilan-Losee}
\setcounter{section}{1}
\begin{document}
\includepdf[pages=-]{homework-5_cover-page.pdf}
\begin{enumerate}[label=\thesubsection]
	\setcounter{section}{11}
	\setcounter{subsection}{2}
	\item The reference information allows us to assume that at full load (at maximum capacity), the given transformer will deliver \(S = 2^{\mathrm{kVA}}\) of apparent power, with a voltage ratio of \(\dfrac{E_1}{E_2}=\dfrac{240^{\mathrm{V}}}{120^{\mathrm{V}}} = 2\).

		\textcolor{red}{An important fact to note about transformers is that the power on the primary and secondary side are \textit{equal}; the differences between the two sides will come down to voltage and current, both of which are a function of the \textit{ratio} between the number of turns in the primary and secondary windings.}

		Assuming all values are given in RMS form,
	\begin{enumerate}[label=\alph*)]
		\item Using the formula for power factor, 
			\[
				\mathrm{P.F.} = \cos{(\theta)} = \frac{P}{S} = 0.8,
			\]
			we may calculate real power \(P\) as 
			\[
				P = 0.8 \times 2^{\mathrm{kVA}} = 1.6^{\mathrm{kW}}.
			\]
			\textcolor{red}{Visualizing the concept of the power triangle, remember that \(\overline{S}\) is the complex power; that is, \(P+jQ\), where \(P\) is real power and \(Q\) is reactive power, and the magnitude of \(|\overline{S}|=S\), or apparent power, is the length of the triangle's hypotenuse, \(P\) is the length of the triangle's base, and \(Q\) is its height (\(Q\) can be negative or positive, determining whether the reactive power is inductive or capacitive, and \(P\) is always positive). Taking the angle between the base and hypotenuse, \(\theta\), we can derive the power factor as \(\mathrm{P.F.} = \cos{(\theta)}\), so that \(P=S\cos{(\theta)}\) and \(Q=S\sin{(\theta)}\).}
		\item Using the formula for RMS current,
			\[
			  I = \frac{S}{E}, 
			\]
			we may calculate load current as 
			\[
				I_{\mathrm{load}} = \frac{2^{\mathrm{kVA}}}{120^{\mathrm{V}}} = 16.67^{\mathrm{A}}.
			\]
			\textcolor{red}{Note that, since the apparent power on the secondary side is equal to the primary side, we can derive the load current by simply moving around the standard power formula, P=IV, except that here we use apparent power rather than just real power. Of course, if we were to use complex power and not just its magnitude, the result would be a complex current value (either in complex or phasor form) with a \textit{magnitude} of \(16.67^{\mathrm{A}}\), so it's important whether the question asks for just the magnitude of the load current (this is usually the case) or the magnitude and phase angle of the load current.}
	\end{enumerate}
	\setcounter{subsection}{5}
	\item
	\begin{enumerate}[label=\alph*)]
			\item For three transformers rated at \(S = 10^{\mathrm{kVA}}\), we may find the total rated power of the transformer bank as
				\[
					S_{\mathrm{total}} = 10^{\mathrm{kVA}} \times 3 = 30^{\mathrm{kVA}}.
				\]
				\textcolor{red}{Regardless of the configuration (delta or wye), mashing three transformers together that each have a rated power of \(S\) just results in \(S \times 3\).}
			\item For a wye configuration, the line-to-line voltage is \(\sqrt{3}\) times greater than the line-to-neutral voltage, so the voltage across \(N_1\) is
				\[
					V_{ab} = 400^{\mathrm{V}} \times \sqrt{3} = 692.82^{\mathrm{V}}.
				\]
				\textcolor{red}{I've played with a couple ideas and I \textit{think} the best way to visualize this is as two, three-phase power source, where each winding is a voltage source, and each transformer is a pair of windings. Using that analogy, we can recycle three-phase power formulas to show that three primary windings, each held at \(400^{\mathrm{V}}\), arranged in a wye configuration, will create three \(400^{\mathrm{V}}\) line-to-neutral voltage sources (think of each winding being connected on one side to the line, and on the other to the neutral, with an equal voltage across each). Again referring back to three-phase power concepts, the line-to-line voltage of a three-phase power source as a function of the line-to-netural voltage is described by \(V_{l} = \sqrt{3}V_{an}\).}


				In the delta configuration, on the other hand, the voltage across \(N_2\) is the same as the line-to-line voltage. 

				\textcolor{red}{Same thing here, except these secondary windings are in \textit{delta} configuration, with each side connected between lines, so that \(V_{LL} = V_{AB}\)}

				Therefore, the line-to-line voltage ratio is
				\[
					\frac{V_{LL_1}}{V_{LL_2}} = \frac{692.82^{\mathrm{V}}}{300^{\mathrm{V}}} = 2.31
				\]
		\end{enumerate}

		\textcolor{red}{Then you just find the ratio between the two. Coincidentally, this is not the same as the turns ratio, since we are not in a wye-wye or delta-delta configuration.}

		\setcounter{subsection}{7}
		\item
			\begin{enumerate}[label=\alph*)]
				\item To find the turns ratio, we may simply use the same assumptions stated in the previous question, such that the turns ratio is
					\[
						\frac{N_1}{N_2} = \frac{(120/\sqrt{3})^{\mathrm{kV}}}{13.8^{\mathrm{kV}}} = 5.02,
					\]
					assuming that this is a \textit{step-down} transformer bank.

					\textcolor{red}{This is an odd one for a couple reasons; the specification given is flipped, since the primary side is usually connected to the transmission line, but oh well, doesn't technically matter other than convention. More importantly, the textbook makes a note that when represented in this format, \(V_1(\Delta)/V_2(\mathrm{Y})\), both voltages are \textbf{line-to-line voltages}, regardless of the arrangement of the transformers (delta or wye). Just like before, we'll convert one side to match the other, you can either convert the line to line voltage to be line to neutral, or vice versa. Above, I converted the line-to-neutral voltage to be line-to-line to maintain consistency, and then simply took the ratio of the two.}

				\item Finding line current as a function of power on the load side, 
					\[
						I = \frac{12^{\mathrm{MVA}}}{(120 \times \sqrt{3})^{\mathrm{kV}}} = 57.74^{\mathrm{A}}. 
					\]

					\textcolor{red}{There are a couple different ways to do this, and the way I did it is not the intuitive way; it was just by using a weird formula I grabbed from the textbook. In actuality, it would make more sense to take the apparent power, \(S\), and divide it by 3 to get the power of a single transformer (since we're only looking for the current on one line \(\to\) one transformer \(\to\) one secondary winding. Then, we may divide that value by the line-to-neutral voltage, found by dividing line-to-line voltage by \(\sqrt{3}\). In other words (numbers),
					\[
						I_L = \frac{S/3}{V_{AB}/\sqrt{3}} = \frac{4^{\mathrm{MVA}}}{69.28^{\mathrm{kV}}} = 57.74^{\mathrm{A}}.
					\]}

			\end{enumerate}
			\setcounter{subsection}{11}
			\item
			\begin{enumerate}[label=\alph*)]
				\item First, we may compute the voltage per turn on the primary side as
					\[
						V_{T} = \frac{E_1}{N_1} = \frac{120^{\mathrm{V}}}{10,000} = 12^{\mathrm{mV}},
					\]

					\textcolor{red}{Voltage per turn is exactly what it sounds like; the total voltage across a winding, divided by the number of turns in the winding.}

					and then compute the voltage on the two secondary windings as 
					\begin{align*}
						&E_2 = V_T \times N_2 = 12^{\mathrm{mV}} \times 5,000 = 60^{\mathrm{V}} \\
						&E_3 = V_T \times N_3 = 12^{\mathrm{mV}} \times 1,000 = 12^{\mathrm{V}}
					\end{align*}
					
					\textcolor{red}{Here, we're just doing the opposite: multiplying the voltage per turn (which should be consistent throughout \textbf{every winding} in the transformer, by the number of turns in that winding to find the total voltage across the winding.}

				\item Computing the current in each winding as a function of the apparent power \(S\) and the winding voltage, 
					\begin{align*}
						&I_1 = \frac{678^{\mathrm{VA}}}{120^{\mathrm{V}}} = 5.64^{\mathrm{A}}\\
						&I_2 = \frac{671^{\mathrm{VA}}}{60^{\mathrm{V}}} = 11.18^{\mathrm{A}}\\
						&I_3 = \frac{43^{\mathrm{VA}}}{12^{\mathrm{V}}} = 3.61^{\mathrm{A}}\\
					\end{align*}
			\end{enumerate}

			\textcolor{red}{Also pretty straightforward, power over voltage gives you current, and power consumed in each winding is provided in the problem description. It's important to note that we we are given real and reactive power separately, so in the above, I took the absolute value to find apparent power, \(S = \sqrt{P^2 + Q^2}\). I did get docked points for not including the phase of each current, though, so it would have been better to use the complex power \(\overline{S}\) in the numerator, and find the corresponding current phasor (though the magnitudes I found are correct).}

			\setcounter{subsection}{12}
		\item Using the equivalence
			\[
			  N_1 I_1 = N_2 I_2 + N_3 I_3,
			\]

			\textcolor{red}{Gonna be honest idk where this formula is derived from, but it seems not terrible to memorize.}

			\begin{align*}
			  &N_2 I_2 = 3000 \times 2^{\mathrm{A}} = 6000 \\
				&N_3 I_3 = 6000 \times 1^{\mathrm{A}} = 6000,
			\end{align*}
			so that \(N_1 I_1 = 12,000\), and
			\[
			  N_1 = \frac{12,000}{I_1} = 1,200.
			\]
\end{enumerate}
\end{document}
